%! AP Statistics — Unit 9, Lesson 9.5 — Guided Notes
\documentclass[10pt]{article}
\usepackage{apstats-article}
\usepackage{apstats-boxes}

\begin{document}

\pageheader{Unit 9, Lesson 9.5}{Guided Notes}
\namedateperiod

\begin{objectivebox}
By the end of this lesson, I will be able to:
\begin{itemize}
    \item[$\square$] \writeline
    \item[$\square$] \writeline
    \item[$\square$] \writeline
    \item[$\square$] \writeline
\end{itemize}
\end{objectivebox}

\begin{vocabbox}
\termblanklong{Null hypothesis for slope ($H_0$: $\beta = 0$)}
\termblanklong{Alternative hypothesis ($H_a$: $\beta \ne 0$)}
\termblanklong{Test statistic ($t = b / SE_b$)}
\termblanklong{$p$-value (in context of regression)}
\termblanklong{Regression output table}
\termblanklong{LINER conditions}
\end{vocabbox}

\begin{hookbox}
A café owner wonders whether outdoor temperature ($x$, in ${}^\circ$F) is linearly
related to daily iced coffee sales ($y$, number of drinks). She collects data for 18
randomly selected days and fits a regression line. Her software produces the following
output:

\begin{center}
\renewcommand{\arraystretch}{1.4}
\begin{tabular}{lrrrr}
    \textbf{Predictor} & \textbf{Coef} & \textbf{SE Coef} & \textbf{T} & \textbf{P} \\
    \hline
    Constant    & 12.34 & 3.21 & 3.84 & 0.002 \\
    Temperature &  2.30 & 0.58 & 3.97 & 0.001 \\
\end{tabular}
\end{center}

What do the four numbers in the \textit{Temperature} row tell us?
\writelines{2}
\end{hookbox}

\begin{notesbox}{Reading a Regression Output Table (VAR-7.I)}
In a regression output table, the \textbf{slope row} (the predictor row, not Constant)
contains the key values for inference.

\begin{center}
\renewcommand{\arraystretch}{1.5}
\begin{tabular}{lcccc}
    \textbf{Predictor} & \textbf{Coef} & \textbf{SE Coef} & \textbf{T} & \textbf{P} \\
    \hline
    Constant   & \blank{2cm} & \blank{2cm} & \blank{1.5cm} & \blank{1.5cm} \\
    \textit{[x variable]} & \blank{2cm} & \blank{2cm} & \blank{1.5cm} & \blank{1.5cm} \\
    \multicolumn{5}{l}{\small\itshape $\uparrow$ read $b$ and $SE_b$ from this row} \\
\end{tabular}
\end{center}

\vspace{0.05in}
\textbf{Café Example} ($n = 18$): $b = $ \blank{1.5cm}, $SE_b = $ \blank{1.5cm},
$t = $ \blank{1.5cm}, $p$-value $= $ \blank{1.5cm}

\vspace{0.1in}
Test statistic formula: \quad $t = \dfrac{b}{\blank{3cm}}$ \qquad
Degrees of freedom: \quad $df = n - \blank{1cm}$
\end{notesbox}

\begin{notesbox}{The 4-Step Significance Test for Slope (VAR-7.I, DAT-3.E/F)}

\textbf{Step 1 — Hypotheses.} Define $\beta$ in words, then state:

\quad H$_0$: $\beta = $ \blank{1.5cm} \quad (no linear relationship)

\quad H$_a$: $\beta \ne $ \blank{1.5cm} \quad (there is a linear relationship)

\vspace{0.12in}
\textbf{Step 2 — Conditions (LINER).}
\begin{itemize}
    \item \textbf{L}inear: \blank{8cm}
    \item \textbf{I}ndependent: \blank{8cm}
    \item \textbf{N}ormal: \blank{8cm}
    \item \textbf{E}qual variance: \blank{8cm}
    \item \textbf{R}andom: \blank{8cm}
\end{itemize}

\vspace{0.12in}
\textbf{Step 3 — Test Statistic \& $p$-value.}

\quad $t = \dfrac{b}{SE_b} = \dfrac{\blank{1.5cm}}{\blank{1.5cm}} = \blank{1.5cm}$
\qquad $df = \blank{1.5cm}$ \qquad $p$\text{-value} $= \blank{2cm}$ (from output)

\vspace{0.12in}
\textbf{Step 4 — Conclusion.}

Since $p$ \blank{1cm} $\alpha$, we \blank{2.5cm} H$_0$.
There \blank{1.5cm} convincing evidence that the true slope $\beta$ \blank{1cm} $0$.

\end{notesbox}

\begin{notesbox}{Worked Example: Iced Coffee Sales (All 4 Steps)}

\textbf{Step 1.} Let $\beta = $ \writeline

\quad H$_0$: \blank{7cm} \quad H$_a$: \blank{7cm}

\vspace{0.1in}
\textbf{Step 2.} LINER conditions: \writelines{3}

\vspace{0.05in}
\textbf{Step 3.} $t = \dfrac{\blank{1.5cm}}{\blank{1.5cm}} = \blank{1.5cm}$
\quad $df = 18 - 2 = \blank{1.5cm}$ \quad $p = \blank{1.5cm}$ (from output)

\vspace{0.1in}
\textbf{Step 4.} Since $p$ \blank{1cm} $0.05$, we \blank{2.5cm} H$_0$.

\writelines{2}

\end{notesbox}

\end{document}
